% ============================================================
% Lecture 3: infinite-width training dynamics
% ============================================================
\chapter{Infinite-Width Training Dynamics}
\label{ch:lecture3}

\begin{keyideas}
\begin{itemize}
    \item The question is how the initialization and learning rate must
    scale with the width $N$ for the limit $N\to\infty$ to exist and stay
    nontrivial. One family covers both answers:
    $f_\theta=\frac{1}{\gamma\sqrt N}\sum_i W^{(2)}_i\phi(W^{(1)}_i\cdot x/\sqrt D)$
    trained with learning rate $\eta\gamma^2$.
    \item \textbf{NTK} ($\gamma=1$): each neuron moves only $O(N^{-1/2})$,
    the kernel $K_\theta$ freezes, and training is kernel regression,
    $\mathbf r(t)=e^{-\eta K_\infty t/n}\,\mathbf r(0)$.
    \item \textbf{Mean field} ($\gamma=\sqrt N$): each neuron moves $O(1)$
    and feels the others only through their empirical distribution $\mu_t$,
    which follows the Wasserstein gradient flow
    $\partial_t\mu_t=\eta\nabla\cdot(\mu_t\nabla\,\delta\mathcal L/\delta\mu_t)$.
    \item Lazy training versus feature learning is decided by the
    parametrization, not by the width alone.
\end{itemize}
\end{keyideas}

\section{Setup}

Consider a two-layer neural network
\begin{equation}
    f_\theta(x)
    =
    s_2
    \sum_{i=1}^{N}
    W_i^{(2)}
    \phi\!\left(
        s_1 W_i^{(1)}\cdot x
    \right),
\end{equation}
where \(x\in\mathbb{R}^D\) is the input, \(N\) is the width, and
\[
\theta_i
=
\left(
W_i^{(1)},W_i^{(2)}
\right)
\]
denotes the parameters associated with neuron \(i\).

The network can be viewed as a bipartite computational graph:
the input is connected to \(N\) hidden neurons through \(W_i^{(1)}\),
and the hidden neurons are connected to the output through \(W_i^{(2)}\).

We train the network by gradient descent,
\begin{equation}
    \theta_{t+1}
    =
    \theta_t
    -
    \eta \nabla_\theta \mathcal{L}(\theta_t),
\end{equation}
with, for example, the squared loss
\begin{equation}
    \mathcal{L}(\theta)
    =
    \frac{1}{2n}
    \sum_{\alpha=1}^{n}
    \left(
        f_\theta(x_\alpha)-y_\alpha
    \right)^2.
\end{equation}

Throughout this discussion, we first keep the number of data points \(n\)
and the training time fixed, and study the limit
\[
N\rightarrow\infty.
\]

A more difficult problem is to study joint limits in which
the number of data points or the training time also grows with \(N\).

\section{Main Question}

The main question is:

\begin{quote}
How should the initialization and learning rate scale with the width \(N\)
so that the limit \(N\rightarrow\infty\) exists and remains nontrivial?
\end{quote}

We take
\begin{equation}
    W_{ij}^{(\ell)}
    \sim
    \mathcal{N}(0,1),
\end{equation}
and introduce the unified parameterization
\begin{equation}
    \boxed{
    f_\theta(x)
    =
    \frac{1}{\gamma\sqrt{N}}
    \sum_{i=1}^{N}
    W_i^{(2)}
    \phi\!\left(
        \frac{1}{\sqrt{D}}
        W_i^{(1)}\cdot x
    \right)
    }
\end{equation}
together with the learning-rate scaling
\begin{equation}
    \boxed{
    \theta_{t+1}
    =
    \theta_t
    -
    \eta\gamma^2
    \nabla_\theta \mathcal{L}(\theta_t).
    }
\end{equation}

Two important regimes are
\begin{align}
    \gamma &= 1
    &&\text{NTK / lazy-training regime},\\
    \gamma &= \sqrt{N}
    &&\text{mean-field regime}.
\end{align}

These two choices lead to qualitatively different infinite-width
training dynamics.

% ============================================================
\section{NTK Regime}

Take
\[
\gamma=1.
\]
Then
\begin{equation}
    f_\theta(x)
    =
    \frac{1}{\sqrt{N}}
    \sum_{i=1}^{N}
    W_i^{(2)}
    \phi\!\left(
        \frac{1}{\sqrt{D}}
        W_i^{(1)}\cdot x
    \right).
\end{equation}

The important phenomenon is that, for large \(N\), the parameters move
only a small distance from their initialization. Therefore the network
can be approximated by its first-order Taylor expansion around
\(\theta_0\):
\begin{equation}
    \boxed{
    f_\theta^{\mathrm{lin}}(x)
    =
    f_{\theta_0}(x)
    +
    \left\langle
        \theta-\theta_0,
        \nabla_\theta f_{\theta_0}(x)
    \right\rangle.
    }
\end{equation}

Schematically, under suitable regularity assumptions and for finite
training time,
\begin{align}
    \sup_{t\leq T}
    \left\|
        f_{\theta_t}
        -
        f_{\theta_t}^{\mathrm{lin}}
    \right\|
    &\longrightarrow 0,
    \\
    \sup_{t\leq T}
    \left\|
        K_{\theta_t}
        -
        K_{\theta_0}
    \right\|
    &\longrightarrow 0
\end{align}
as \(N\rightarrow\infty\).

Thus, in the NTK limit, training of the neural network becomes
equivalent to training a linear model with the fixed feature map
\[
x
\mapsto
\nabla_\theta f_{\theta_0}(x).
\]

In this sense,

\begin{quote}
At infinite width, the nonlinear neural network has linearized training
dynamics around its initialization.
\end{quote}

\section{Derivation of NTK Dynamics}

It is convenient to consider continuous-time gradient flow,
\begin{equation}
    \frac{d\theta_t}{dt}
    =
    -\eta
    \nabla_\theta
    \mathcal{L}(\theta_t).
\end{equation}

Then
\begin{align}
    \frac{d}{dt}
    f_{\theta_t}(x)
    &=
    \nabla_\theta f_{\theta_t}(x)
    \cdot
    \frac{d\theta_t}{dt}
    \\
    &=
    -\eta
    \nabla_\theta f_{\theta_t}(x)
    \cdot
    \nabla_\theta
    \mathcal{L}(\theta_t).
\end{align}

For the squared loss,
\begin{equation}
    \nabla_\theta\mathcal{L}
    =
    \frac{1}{n}
    \sum_{\alpha=1}^{n}
    \left(
        f_{\theta}(x_\alpha)-y_\alpha
    \right)
    \nabla_\theta f_\theta(x_\alpha).
\end{equation}

Therefore,
\begin{equation}
    \boxed{
    \frac{d}{dt}
    f_{\theta_t}(x)
    =
    -
    \frac{\eta}{n}
    \sum_{\alpha=1}^{n}
    \left(
        f_{\theta_t}(x_\alpha)-y_\alpha
    \right)
    K_{\theta_t}(x,x_\alpha),
    }
\end{equation}
where
\begin{equation}
    \boxed{
    K_\theta(x,x')
    =
    \left\langle
        \nabla_\theta f_\theta(x),
        \nabla_\theta f_\theta(x')
    \right\rangle
    }
\end{equation}
is the \textbf{Neural Tangent Kernel (NTK)}.

For the training data, define
\[
f_\alpha(t)=f_{\theta_t}(x_\alpha),
\qquad
r_\alpha(t)=f_\alpha(t)-y_\alpha.
\]
Then
\begin{equation}
    \boxed{
    \dot{\mathbf r}(t)
    =
    -
    \frac{\eta}{n}
    K_{\theta_t}
    \mathbf r(t).
    }
\end{equation}

If the kernel becomes constant,
\[
K_{\theta_t}
\approx
K_{\theta_0}
\rightarrow
K_\infty,
\]
then
\begin{equation}
    \boxed{
    \mathbf r(t)
    =
    \exp\!\left(
        -\frac{\eta}{n}K_\infty t
    \right)
    \mathbf r(0).
    }
\end{equation}

Thus the entire training dynamics is determined by a fixed kernel.

\section{Why Does the NTK Freeze?}

Because the network is a sum over neurons,
\begin{equation}
    K_\theta(x,x')
    =
    \frac{1}{N}
    \sum_{i=1}^{N}
    K_{\theta_i}(x,x'),
\end{equation}
so \(K_\theta=O(1)\).

For a two-layer network, the Hessian with respect to neuron parameters
has a block-diagonal structure,
\begin{equation}
    \operatorname{Hess}_\theta f
    =
    \operatorname{diag}
    \left[
        \operatorname{Hess}_{\theta_1}f,
        \ldots,
        \operatorname{Hess}_{\theta_N}f
    \right].
\end{equation}

In NTK parameterization, each neuron contributes only
\(O(N^{-1/2})\) to the network output. Correspondingly, each individual
neuron moves only
\[
\Delta\theta_i
=
O(N^{-1/2})
\]
over an \(O(1)\) training time.

Hence
\[
\theta_i(t)
\approx
\theta_i(0),
\]
and the features generated by each neuron barely change.

Therefore
\begin{equation}
    \boxed{
    \frac{d}{dt}K_{\theta_t}
    \longrightarrow 0
    \qquad
    \text{as }
    N\rightarrow\infty.
    }
\end{equation}

The nonlinear terms in the Taylor expansion vanish, leaving only the
linearized dynamics.

\[
\boxed{
\text{NTK limit}
\quad\Longrightarrow\quad
\text{kernel learning with nearly frozen features}.
}
\]

% ============================================================
\section{Mean-Field Regime}

Now take
\[
\gamma=\sqrt{N}.
\]
Then
\begin{equation}
    \boxed{
    f_\theta(x)
    =
    \frac{1}{N}
    \sum_{i=1}^{N}
    W_i^{(2)}
    \phi\!\left(
        \frac{1}{\sqrt{D}}
        W_i^{(1)}\cdot x
    \right).
    }
\end{equation}

The learning-rate scaling becomes
\begin{equation}
    \boxed{
    \theta_{t+1}
    =
    \theta_t
    -
    \eta N
    \nabla_\theta
    \mathcal{L}(\theta_t).
    }
\end{equation}

Unlike the NTK regime, individual neurons now move by an \(O(1)\)
amount during training. Therefore the feature representation can evolve
nontrivially.

Define
\begin{equation}
    \vartheta_i
    =
    \left(
        W_i^{(1)},W_i^{(2)}
    \right)
\end{equation}
and the empirical distribution of neurons
\begin{equation}
    \boxed{
    \mu_t^N
    =
    \frac{1}{N}
    \sum_{i=1}^{N}
    \delta_{\vartheta_i(t)}.
    }
\end{equation}

Then the network output can be written as
\begin{equation}
    f_{\mu_t^N}(x)
    =
    \int
    W^{(2)}
    \phi\!\left(
        \frac{1}{\sqrt{D}}
        W^{(1)}\cdot x
    \right)
    d\mu_t^N(W^{(1)},W^{(2)}).
\end{equation}

As
\[
N\rightarrow\infty,
\]
the discrete empirical measure \(\mu_t^N\) approaches a continuous
probability measure \(\mu_t\).

Thus instead of describing \(N\) individual neurons, we describe the
evolution of the distribution of neurons.

\section{Dynamics of an Individual Neuron}

Define
\begin{equation}
    \psi(\vartheta;x)
    =
    W^{(2)}
    \phi\!\left(
        \frac{1}{\sqrt{D}}
        W^{(1)}\cdot x
    \right).
\end{equation}

Then
\[
f_{\mu}(x)
=
\int
\psi(\vartheta;x)
\,d\mu(\vartheta).
\]

The gradient of the network with respect to neuron \(i\) is
\begin{equation}
    \nabla_{\vartheta_i}
    f_\theta(x_\alpha)
    =
    \frac{1}{N}
    \nabla_{\vartheta_i}
    \psi(\vartheta_i;x_\alpha).
\end{equation}

Because the learning rate is scaled by \(N\),
\begin{align}
    \frac{d}{dt}
    \vartheta_i(t)
    &=
    -\eta N
    \nabla_{\vartheta_i}
    \mathcal{L}
    \\
    &=
    -\frac{\eta}{n}
    \sum_{\alpha=1}^{n}
    \left(
        f_{\mu_t}(x_\alpha)-y_\alpha
    \right)
    \nabla_{\vartheta_i}
    \psi(\vartheta_i;x_\alpha).
\end{align}

Therefore,
\begin{equation}
    \boxed{
    \dot{\vartheta}_i
    =
    -
    \eta
    \mathcal{M}
    (\vartheta_i,\mu_t).
    }
\end{equation}
with
\begin{equation}
    \mathcal{M}(\vartheta,\mu)
    =
    \frac{1}{n}
    \sum_{\alpha=1}^{n}
    \left(
        f_\mu(x_\alpha)-y_\alpha
    \right)
    \nabla_\vartheta
    \psi(\vartheta;x_\alpha),
\end{equation}
which the next section identifies as
\(\nabla_\vartheta\,\delta\mathcal{L}/\delta\mu\).

The important point is that neuron \(i\) does not depend directly on
all other neurons separately. It depends on them only through the
empirical distribution
\[
\mu_t.
\]

This is a mean-field interaction.

\section{Evolution of the Neuron Distribution}

The loss can now be regarded as a functional of the measure,
\begin{equation}
    \mathcal{L}[\mu]
    =
    \frac{1}{2n}
    \sum_{\alpha=1}^{n}
    \left(
        f_\mu(x_\alpha)-y_\alpha
    \right)^2.
\end{equation}

Its functional derivative is
\begin{equation}
    \frac{\delta\mathcal{L}}{\delta\mu}
    (\vartheta)
    =
    \frac{1}{n}
    \sum_{\alpha=1}^{n}
    \left(
        f_\mu(x_\alpha)-y_\alpha
    \right)
    \psi(\vartheta;x_\alpha).
\end{equation}

Therefore the particle dynamics becomes
\begin{equation}
    \boxed{
    \dot{\vartheta}
    =
    -
    \eta
    \nabla_\vartheta
    \frac{\delta\mathcal{L}}
         {\delta\mu}(\vartheta).
    }
\end{equation}

If the particles follow a velocity field
\[
v_t(\vartheta)
=
-\eta
\nabla_\vartheta
\frac{\delta\mathcal{L}}
     {\delta\mu_t}(\vartheta),
\]
conservation of probability gives the continuity equation
\begin{equation}
    \partial_t\mu_t
    +
    \nabla_\vartheta\cdot
    \left(
        \mu_t v_t
    \right)
    =
    0.
\end{equation}

Substituting \(v_t\),
\begin{equation}
    \boxed{
    \partial_t\mu_t
    =
    \eta
    \nabla_\vartheta\cdot
    \left(
        \mu_t
        \nabla_\vartheta
        \frac{\delta\mathcal{L}}
             {\delta\mu_t}
    \right).
    }
\end{equation}

This is a \textbf{Wasserstein gradient flow} of the loss functional
\(\mathcal{L}[\mu]\).

Indeed,
\begin{equation}
    \frac{d}{dt}\mathcal{L}[\mu_t]
    =
    -
    \eta
    \int
    \left\|
        \nabla_\vartheta
        \frac{\delta\mathcal{L}}
             {\delta\mu_t}
    \right\|^2
    d\mu_t(\vartheta)
    \leq 0.
\end{equation}

Thus gradient descent on a finite collection of neurons becomes,
at infinite width, gradient flow of a probability distribution in
parameter space.

\section{NTK vs.\ Mean Field}

The two limits therefore have very different physical pictures:
\begin{align}
    \text{NTK:}\qquad
    &\theta_i(t)-\theta_i(0)
    \rightarrow 0,
    \\
    &K_{\theta_t}
    \approx
    K_{\theta_0},
    \\
    &\text{features are approximately frozen};
    \\[1em]
    \text{Mean Field:}\qquad
    &\theta_i(t)-\theta_i(0)
    =
    O(1),
    \\
    &\mu_0
    \rightarrow
    \mu_t,
    \\
    &\text{the neuron distribution and features evolve}.
\end{align}

Hence,
\[
\boxed{
\begin{array}{c}
\text{NTK limit}
\\[2mm]
\text{fixed kernel / lazy feature learning}
\end{array}
\qquad\text{vs.}\qquad
\begin{array}{c}
\text{Mean-field limit}
\\[2mm]
\text{distributional dynamics / feature learning}.
\end{array}
}
\]

The central conceptual change in the mean-field limit is
\[
\boxed{
\text{many interacting neurons}
\quad\longrightarrow\quad
\text{a probability measure evolving by optimal transport}.
}
\]

% WEBFIG: widthdynamics
